\RequirePackage{plautopatch}
\documentclass[14pt,aspectratio=169,xcolor=dvipsnames,table,onlytextwidth,dvipdfmx]{beamer}
\usepackage[utf8]{inputenc}
\usepackage{bxdpx-beamer} % dvipdfmxなので必要

\usetheme{Boadilla}
%Beamerフォント設定
\usefonttheme{professionalfonts}
\usepackage[T1]{fontenc}
\usepackage[deluxe,uplatex]{otf} % 日本語多ウェイト化
\usepackage{mlmodern}  % 太いComputer Modern
% MLmodernのバグを修正: cf. https://tex.stackexchange.com/questions/646333/size-of-integral-symbol-in-section-header-with-mlmodern
\DeclareFontFamily{OMX}{mlmex}{}
\DeclareFontShape{OMX}{mlmex}{m}{n}{%
   <->mlmex10%
   }{}
\usepackage{newtxtext} % 数式以外をTXフォントで上書き
% \usepackage[noalphabet,unicode]{pxchfon}
% \setminchofont{KozMinPro-Regular.otf}% 游明朝Regular
% \setboldminchofont{KozMinPro-Bold.otf}% 游明朝Demibold
% \setgothicfont[0]{BIZ-UDGothicR.ttc}% BIZ UDゴシックR
% \setboldgothicfont[0]{BIZ-UDGothicB.ttc}% BIZ UDゴシックB
\renewcommand{\familydefault}{\sfdefault}  % 英文をサンセリフ体に
\renewcommand{\kanjifamilydefault}{\gtdefault}  % 日本語をゴシック体に
\usefonttheme{structurebold} % タイトル部を太字
\setbeamerfont{alerted text}{series=\bfseries} % Alertを太字
\setbeamerfont{section in toc}{series=\mdseries} % 目次は太字にしない
\setbeamerfont{frametitle}{size=\Large} % フレームタイトル文字サイズ
\setbeamerfont{title}{size=\LARGE} % タイトル文字サイズ
\setbeamerfont{date}{size=\small}  % 日付文字サイズ
\usepackage{bxcoloremoji}

% Babel (日本語の場合のみ・英語の場合は不要)
\uselanguage{japanese}
\languagepath{japanese}
\deftranslation[to=japanese]{Theorem}{定理}
\deftranslation[to=japanese]{Lemma}{補題}
\deftranslation[to=japanese]{Example}{例}
\deftranslation[to=japanese]{Examples}{例}
\deftranslation[to=japanese]{Definition}{定義}
\deftranslation[to=japanese]{Definitions}{定義}
\deftranslation[to=japanese]{Problem}{問題}
\deftranslation[to=japanese]{Solution}{解}
\deftranslation[to=japanese]{Fact}{事実}
\deftranslation[to=japanese]{Proof}{証明}
\deftranslation[to=japanese]{Corollary}{系}
\setbeamertemplate{theorems}[normal font] %日本語の場合，定理内を斜体にしない
% overlay-aware proof environment
\renewenvironment<>{proof}{%
\par
\begin{actionenv}#1%
\structure{(証明)}}
{\qed\end{actionenv}}

%Beamer色設定
\definecolor{UniBlue}{RGB}{0,150,200}
\definecolor{UniGreen}{RGB}{3,175,122}
\definecolor{UniPink}{RGB}{255,128,255}
\definecolor{AlertOrange}{RGB}{255,76,0}
\definecolor{AlmostBlack}{RGB}{38,38,38}
\setbeamercolor{normal text}{fg=AlmostBlack}  % 本文カラー
\setbeamercolor{structure}{fg=UniBlue} % 見出しカラー
\setbeamercolor{block title}{fg=UniBlue!50!black} % ブロック部分タイトルカラー
\setbeamercolor{alerted text}{fg=AlertOrange} % \alert 文字カラー
\setbeamercolor{bibliography entry author}{fg=AlmostBlack} % Bib 文字カラー
\setbeamercolor{bibliography entry note}{fg=AlmostBlack} % Bib 文字カラー
\mode<beamer>{
    \definecolor{BackGroundGray}{RGB}{254,254,254}
    \setbeamercolor{background canvas}{bg=BackGroundGray} % スライドモードのみ背景をわずかにグレーにする
}

%フラットデザイン化
\setbeamertemplate{blocks}[rounded] % Blockの影を消す
\useinnertheme{circles} % 箇条書きをシンプルに
\setbeamertemplate{navigation symbols}{} % ナビゲーションシンボルを消す
\setbeamertemplate{footline}[frame number] % フッターはスライド番号のみ

%タイトルページ
\setbeamertemplate{title page}{%
    \vspace{2.5em}
    {\usebeamerfont{title} \usebeamercolor[fg]{title} \inserttitle \par}
    {\usebeamerfont{subtitle}\usebeamercolor[fg]{subtitle}\insertsubtitle \par}

    \vspace{1.5em}
    \begin{columns}[]
        \begin{column}{.35\linewidth}
        \usebeamerfont{titlegraphic}\inserttitlegraphic\par
        \end{column}
        \begin{column}{.6\linewidth}\setlength\topsep{0pt}
            \begin{flushright}
        \usebeamerfont{author}\insertauthor\par
        \usebeamerfont{institute}\insertinstitute \par
        \vspace{3em}
        \usebeamerfont{date}\insertdate\par
            \end{flushright}
        \end{column}
    \end{columns}
}

%biblatex
% \usepackage[style=ext-authoryear,citestyle=authoryear,natbib=true,giveninits=true,maxbibnames=99,maxcitenames=10,url=false,isbn=false,doi=false,articlein=false]{biblatex}
% \DeclareNameAlias{author}{first-last}
% \addbibresource{shrunk.bib}
% \newcommand{\mkbibbracketscol}[1]{\textcolor{UniBlue}{\scriptsize \mkbibbrackets{#1}}}
% \DeclareCiteCommand{\cite}[\mkbibbracketscol]
%   {\usebibmacro{prenote}}
%   {\usebibmacro{citeindex}%
%    \usebibmacro{cite}}
%   {\multicitedelim}
%   {\usebibmacro{postnote}}
% % \renewcommand{\cite}[1]{\textcolor{UniBlue}{\scriptsize\citep{#1}}}
% \renewcommand*{\finalnamedelim}{\addcomma\addspace} % remove "and" before the last author
% \DeclareDelimFormat{nameyeardelim}{\addspace}
% \setbeamertemplate{bibliography item}{}

% Algorithm系
\usepackage{algorithm}
\usepackage[noend]{algorithmic}
\algsetup{linenosize=\color{fg!50}\footnotesize}
\renewcommand\algorithmicdo{:}
\renewcommand\algorithmicthen{:}
\renewcommand\algorithmicrequire{\textbf{Input:}}
\renewcommand\algorithmicensure{\textbf{Output:}}
\renewcommand\algorithmiccomment[1]{\hfill\textcolor{fg!50}{\footnotesize\textit{// #1}}}

% TikZ
\usepackage{tikz}
\usetikzlibrary{positioning,shapes,arrows,quotes,shapes.callouts,calc,graphs,graphs.standard,fit,backgrounds,perspective,matrix}
\tikzset{point/.style={circle, fill=AlertOrange, inner sep=1pt, text width=1pt}}
\tikzset{bpoint/.style={circle, fill=black, inner sep=1pt, text width=1pt}}
\tikzset{alt/.code args={<#1>#2#3}{\alt<#1>{\pgfkeysalso{#2}}{\pgfkeysalso{#3}}}}
\tikzset{>=latex}
\tikzset{mynodes/.style={circle,white,fill=UniBlue!90,text width=.8em,inner sep=0pt,text centered,font=\footnotesize}}
\tikzset{myedges/.style={thick}}
\tikzset{mypath/.style={ultra thick, draw=#1}, mypath/.default=AlertOrange}
\tikzset{mysubset/.style={draw,#1,dotted,thick,fill=#1!10}, mysubset/.default=UniBlue}
\tikzset{mypath2/.style={ultra thick, dashed, draw=UniGreen}}
\tikzset{graphs/mygraph/.style={nodes={mynodes}, edges={myedges}, empty nodes}}
\tikzset{nomargin/.style={inner sep=0pt, outer sep=0pt}}
%% graph macros
\tikzgraphsset{declare={mybipartite}{%
subgraph I_nm [V={1,2,3,4,5,6}, W={a,b,c,d,e,f}];
1 -- {a,b};
2 -- {b,c,d,e};
3 -- {d,f};
4 -- {e,f};
{5,6} -- f;
}}
\tikzgraphsset{declare={mybipartitedi}{%
subgraph I_nm [V={1,2,3,4,5,6}, W={a,b,c,d,e,f}];
1 -> {a,b};
2 -> {b,c,d,e};
3 -> {d,f};
4 -> {e,f};
{5,6} -> f;
}}
\tikzgraphsset{declare={mybipartite2}{%
subgraph I_nm [V={1,2,3}, W={a,b,c}];
1 -- {a,b}; 2 -- {a,b,c}; 3 -- {a,c};
}} 
% Petersen graph
\tikzgraphsset{declare={Petersen}{%
    subgraph C_n [n=5,name=A, radius=1.5cm]; 
    subgraph I_n [n=5,name=B, radius=.8cm];
    \foreach \i [evaluate={\j=int(mod(\i+2,5)+1);}] in {1,...,5}{
        A \i -- B \i;
        B \i -- B \j;
}}}
% Example graph for Menger's theorem
\tikzgraphsset{declare={menger}{%
    s[at={(-1,0)}, as={$s$}];
    (s) -> 1[at={(0,1)}] -> 2[at={(1,0)}];
    (s) -> 3[at={(0,-1)}] -> 2;
    4[at={(2,1)}] <- 5[at={(2,-1)}];
    {5,6[at={(3,1)}], 7[at={(3,-1)}]} -> t[at={(4,0)}, as={$t$}];
    (s) -> (2);
    (2) -> {(4), (5)};
    {(4),(5)} -> {(6),(7)};
    (6) -> (7);
    1 -> 4;
}} 
\tikzgraphsset{declare={undirected menger}{%
    s[at={(-1,0)}, as={$s$}];
    (s) -- 1[at={(0,1)}] -- 2[at={(1,0)}];
    (s) -- 3[at={(0,-1)}] -- 2;
    4[at={(2,1)}] -- 5[at={(2,-1)}];
    {5,6[at={(3,1)}], 7[at={(3,-1)}]} -- t[at={(4,0)}, as={$t$}];
    (s) -- (2);
    (2) -- {(4), (5)};
    {(4),(5)} -- {(6),(7)};
    (6) -- (7);
    1 -- 4;
}} 

% Appendix
\usepackage{appendixnumberbeamer}

% macros
\newcommand{\R}{\mathbb{R}}
\newcommand{\Z}{\mathbb{Z}}
\newcommand{\be}{\mathbf{e}}
\newcommand{\zeros}{\mathbf{0}}
\newcommand{\ones}{\mathbf{1}}
\newcommand{\bfzero}{\zeros}
\newcommand{\bfone}{\ones}

% \newcommand{\bp}{\mathbf{p}}
\newcommand{\eps}{\varepsilon}
\newcommand{\defiff}{\overset{\text{def}}{\iff}}
\DeclareMathOperator{\poly}{poly}
\DeclareMathOperator{\conv}{conv}
\DeclareMathOperator{\In}{In}
\DeclareMathOperator{\Out}{Out}
\DeclareMathOperator{\val}{val}
\newcommand{\lo}{{\mathrm{lo}}}
\newcommand{\up}{{\mathrm{up}}}
\newcommand{\caI}{\mathcal{I}}
\newcommand{\caB}{\mathcal{B}}
\newcommand{\caC}{\mathcal{C}}
\newcommand{\bM}{\mathbf{M}}
\newcommand{\IO}{\mathrm{IO}}

\usepackage{mathtools}
\DeclarePairedDelimiter{\abs}{\lvert}{\rvert}
\DeclarePairedDelimiter{\norm}{\lVert}{\rVert}
\DeclarePairedDelimiter{\inprod}{\langle}{\rangle}
\usepackage[no-test-for-array]{nicematrix}
\usepackage{diagbox}

\newcommand{\highlightcap}[3][red]{\tikz[baseline=(x.base)]{\node[rectangle,rounded corners,fill=#1!10](x){#2} node[below=0pt of x, color=#1]{#3};}}
\newcommand{\highlight}[2][red]{\tikz[baseline=(x.base)]{\node[rectangle,rounded corners,fill=#1!10](x){#2};}}
\newcommand{\mycite}[2][\footnotesize]{\textcolor{UniBlue}{#1 [#2]}}

%section
% \AtBeginSection[]{
%     \frame[c]{
%         \centering
%         \begin{tikzpicture}
%             \node[white, circle, fill=UniBlue, text centered, text width=.3\linewidth](bullet) {
%                 {\Huge \bf \insertsectionnumber} \\[4em]
%     };
%     \node[text width=.65\linewidth, right=1em of bullet]{%
%         \usebeamerfont{title}\usebeamercolor[fg]{title}\insertsection\par %
%     };
%         \end{tikzpicture}
%     } %目次スライド
% }

%grid
% \usepackage[colorgrid,gridunit=pt,texcoord]{eso-pic}
\newcommand{\postit}[1]{\tikz[overlay, remember picture]{\node[draw,anchor=north east, align=center, rectangle, fill=yellow!10, xshift=-1pt, yshift=-1pt] at (current page.north east) {#1};}}

\usepackage{tcolorbox}
\newtcolorbox{primal}{colback=red!10!white, colframe=red, title=主問題 (P), left=2pt,right=2pt,top=0pt,bottom=0pt}
\newtcolorbox{dual}{colback=blue!10!white, colframe=blue, title=双対問題 (D), left=2pt,right=2pt,top=0pt,bottom=0pt}

% visible on
\usetikzlibrary{overlay-beamer-styles}
\usepackage[normalem]{ulem}
\usepackage{pxrubrica}
\usepackage{booktabs}

\newcommand{\Aphiup}{\textcolor{UniBlue}{A_\varphi^\up}}
\newcommand{\Aphilow}{\textcolor{AlertOrange}{A_\varphi^\lo}}
\newcommand{\bv}{\boldsymbol{v}}
\newcommand{\bx}{\boldsymbol{x}}
\newcommand{\by}{\boldsymbol{y}}
\newcommand{\bz}{\boldsymbol{z}}
\newcommand{\bw}{\boldsymbol{w}}
\newcommand{\bc}{{\boldsymbol c}}
\DeclareMathOperator{\supp}{supp}
\DeclareMathOperator{\argmax}{argmax}
\DeclareMathOperator{\argmin}{argmin}
\DeclareMathOperator{\E}{\mathbb{E}}
\csname endofdump\endcsname
\title{組合せ最適化特論}
\subtitle{第3回（二部マッチング②）}
\author{担当: \textbf{相馬 輔}}
\date{2025/11/6}

\begin{document}
\maketitle

\begin{frame}{}
    \centering\small
    \tikzset{block/.style={rectangle, draw=gray, fill=gray!10, thick, minimum width=3em, minimum height=2em, align=center}}
    \begin{tikzpicture}
        \node[block,] (BM) {重みなし\\二部マッチング};
        \node[block,above=1em of BM.north west, xshift=-2em] (WBM) {重みつき\\二部マッチング};
        \node[block,above=5em of BM.north east, xshift=1em] (NF) {最大流};
        \node[block,above=9em of BM] (MCF) {最小費用流};
        % \node[block,above=1em of NF.north east, xshift=1em] (MC) {最小カット};
        \node[block,right=11em of BM] (MTR) {マトロイド};
        \node[block,above=7em of MTR] (MI) {マトロイド交差};
        \node[block,above=3em of MCF, xshift=7em] (SFM) {劣モジュラ最小化};
        \graph[mygraph]{
            (BM) -> {(WBM), (NF)} -> (MCF);
            {(MTR), (BM)} -> (MI);
            % (NF) -> (MC);
            {(MCF), (MI)} -> (SFM);
        };
        \begin{scope}[on background layer]
            \node[mysubset, fit=(BM) (WBM), "二部マッチング" {below, UniBlue}]{};
            \node[mysubset, AlertOrange, fill=AlertOrange!10, fit=(NF) (MCF), "ネットワークフロー" {above left, AlertOrange}]{};
            \node[mysubset, UniGreen, fill=UniGreen!10, fit=(MTR) (MI), "マトロイド" {below, UniGreen}]{};
        \end{scope}
    \end{tikzpicture}
\end{frame}

\begin{frame}[allowframebreaks]
    \frametitle{各回の内容（予定）} 
    \small
    \begin{enumerate}
        \item (4/23) イントロ+多面体的組合せ論（線形計画法の復習，整数多面体，完全単模行列）
        \item[] \textcolor{gray}{(4/30) 休み}
        \item (5/7) 二部マッチング①（Konig-Egervaryの定理，増加道アルゴリズム，ハンガリー法）
        \item (5/14) 二部マッチング②（最短路問題の復習，逐次最短路法と主双対法，最適性基準からの見方）
        \item (5/21) 最大流① （定式化，最大流最小カット定理，応用例）
        \item (5/28) 最大流②（残余ネットワーク，Ford--Fulkerson法，Edmonds--Karpのアルゴリズム）+ 最小費用流①（定式化，輸送問題，最大流との関係）
        \item[] \textcolor{gray}{(6/4) 休み}
        % \item[] \sout{(12/11) 最小カット（永持茨木のアルゴリズム，Kargerのアルゴリズム）}
        \item (6/11) 最小費用流②（逐次最短路法，容量スケーリング法）
        \item[] \textcolor{gray}{(6/18) 休み}
        \item[] \textcolor{gray}{(6/25) 休み}
        \item (7/2) マトロイド（定義と公理系，貪欲法，マトロイド多面体） 
        \item (7/9) 独立マッチング・マトロイド交差 （定義，応用例，Edmondsの最大最小定理，増加道アルゴリズム）
        \item (7/16) 劣モジュラ関数①（諸例，劣モジュラ基多面体，Lov\'asz拡張，劣モジュラ最小化）
        \item[] \textcolor{gray}{(7/23) 休み}
        \item (7/30) 劣モジュラ関数②（劣モジュラ最大化，近似アルゴリズム，貪欲法）
        \item (どこか) 精選トピック①
        \item (どこか) 精選トピック②
        \item (どこか) 精選トピック③
    \end{enumerate}

    \bigskip
    \begin{itemize}
        \item レポート課題を6月と期末に出す予定
        \item 精選トピックの日時は今後調整 
    \end{itemize}
\end{frame}

\begin{frame}<1>[label=outline]{目次}
\begin{tikzpicture}
    \tikzstyle{block} = [rectangle, text width=\textwidth, rounded corners, minimum height=4em]
    \tikzstyle{line} = [draw, -latex]

    \node[block,alt={<1,3->{fill=cyan!10}{fill=red!10}}] (combopt) {
        \begin{columns}
            \begin{column}{0.5\linewidth}
        \textbf{1. 最短路問題}
            \end{column}
            \begin{column}{0.4\linewidth}
                \small
                \structure{・} 定義\\
                \structure{・} 負閉路とポテンシャル
            \end{column}
        \end{columns}
    };
    \node[block,alt={<1-2,4->{fill=cyan!10}{fill=red!10}}, below=.5em of combopt] (LP) {
        \begin{columns}
            \begin{column}{0.5\linewidth}
        \textbf{2. 逐次最短路法と主双対法}
            \end{column}
            \begin{column}{0.4\linewidth}
                \small
                \structure{・} 逐次最短路法 \\
                \structure{・} 主双対法
            \end{column}
        \end{columns}
    };
    \node[block,alt={<1-3,5->{fill=cyan!10}{fill=red!10}}, below=.5em of LP] (TU) {
        \begin{columns}
            \begin{column}{0.5\linewidth}
        \textbf{3. 最適性条件からの見方}
            \end{column}
            \begin{column}{0.4\linewidth}
                \small
                \structure{・} ポテンシャル最適性条件 \\
                \structure{・} 閉路最適性条件 \\
                \structure{・} アルゴリズム再訪
            \end{column}
        \end{columns}
    };
\end{tikzpicture}%
\end{frame}

%%%%%%%%%%%%%%%%%%%%%%%%%%%%%%%%%%%%%%%%%%%%%
\section{最短路問題}
\againframe<2|handout:0>{outline}
%%%%%%%%%%%%%%%%%%%%%%%%%%%%%%%%%%%%%%%%%%%%%
\begin{frame}{最短路問題}
     $G = (V, A)$: 有向グラフ, $s, t \in V$, $\ell: A \to \R$: 枝長

     \bigskip
     $P$: $s$--$t$路（$s$から$t$へ行く経路で，同じ頂点を複数回通っても良い） に対し，
     \[
        \ell(P) := \sum_{e \in E(P)} \ell(e)
     \]
     を$P$の長さとする．
     
     \bigskip
     \begin{block}{最短路問題}
        $s$--$t$路$P$のうち，長さ$\ell(P)$最小のものを求めよ．
     \end{block} 
\end{frame}

\begin{frame}{最短路問題と負閉路}
    $s$を固定する．任意の$t$に対して，最短路は存在するか？

    % \bigskip
    % \structure{存在しない場合:}
    % \begin{itemize}
    %     \item $s$から到達不可能な$t$が存在する場合
    %     \item 
    % \end{itemize}
   
    \bigskip
    \begin{lemma}
        $G$の任意の頂点は$s$から到達可能であるとする．このとき，以下は同値:  
        \begin{itemize}
            \item 任意の$t$に対し$s$--$t$最短路が存在 
            \item 有向閉路$C$で$\ell(C) < 0$なるもの（\alert{負閉路}）は存在しない．
        \end{itemize}
    \end{lemma}
\end{frame}

\begin{frame}{ポテンシャル}
    \begin{definition}[ポテンシャル]
        $p: V \to \R$が枝長$\ell$に関する\alert{（実行可能）ポテンシャル} \\ 
        $\defiff$ $p(j) - p(i) \leq \ell(a)$ \quad ($a = ij \in A$)
    \end{definition}

    \begin{lemma}
        $G$の任意の頂点は$s$から到達可能とする．このとき，以下は同値: 
        \begin{enumerate}
            \item 任意の$t$に対し$s$--$t$最短路が存在
            \item ポテンシャル$p: V \to \R$が存在
            \item 負閉路が存在しない
        \end{enumerate}
    \end{lemma}
\end{frame}

\begin{frame}{ポテンシャル}
    \small
    \postit{
        \footnotesize
        \begin{minipage}{.42\textwidth}
        \begin{enumerate}
            \item 任意の$t$に対し$s$--$t$最短路が存在
            \item ポテンシャル$p: V \to \R$が存在
            \item 負閉路が存在しない
        \end{enumerate}
        \end{minipage}
    }

    \bigskip
    \structure{(証明)} 

    \textbf{①$\implies$②:} $p(t) :=$ ($s$--$t$最短路長)とおく．任意の枝$a = ij \in A$
    に対し，図より
    \[ p(j) \leq p(i) + \ell(a). \]

    \bigskip\pause
    \textbf{②$\implies$③:} 任意の有向閉路$C: v_0;v_1;\dots;v_k = v_0$に対して，
    \[ 
        \ell(C) = \sum_{i=1}^k \ell(v_{i-1}v_i) \geq \sum_{i=1}^k (p(v_i) - p(v_{i-1})) = 0. 
    \]

    \bigskip\pause
    \textbf{③$\implies$①:} すでにやった． \qed
\end{frame}

\begin{frame}{最短路問題のアルゴリズム}
    単一の始点$s$に対して，全ての$t$に関する最短路問題を同時に解くアルゴリズムが知られている:  

    \bigskip
    \begin{description}
        \setlength{\itemsep}{1em}
        \item[一般の枝長$\ell$:] \alert{Bellman--Ford法} $O(mn)$時間  {\hfill \footnotesize ※負閉路がある場合は検出も可能}
        \item[非負の枝長$\ell$:] \alert{Dijkstra法} $O(m + n \log n)$時間
    \end{description}
    {\footnotesize\hfill ($n = |V|, m = |E|$)}

    \bigskip
    これらは，$s$--$t$最短路だけでなく，実行可能ポテンシャル(=最短路長)も同時に求める．
\end{frame}

%%%%%%%%%%%%%%%%%%%%%%%%%%%%%%%%%%%%%%%%%%%%%
\section{逐次最短路法と主双対法}
\againframe<3|handout:0>{outline}
%%%%%%%%%%%%%%%%%%%%%%%%%%%%%%%%%%%%%%%%%%%%%
\begin{frame}[label=w-bipartite]
    \frametitle{重み付き二部マッチング}

    \begin{block}{重み付き二部マッチング問題}
        \structure{入力} $G=(V; E)$: 二部グラフ，枝重み$w_e$ ($e \in E$) \\
        \structure{出力} $w(M) := \sum_{e\in M}w_e$が\textbf{最小}のマッチング$M$
    \end{block}
\end{frame}


\begin{frame}
    \frametitle{補助ネットワーク（復習）}

    \small
    \begin{columns}[T]
    \begin{column}{.7\textwidth}
    マッチング$M$に対し，$G$の枝を
    \begin{itemize}
        \item $M$の枝は両向き，
        \item $E\setminus M$の枝は$V^+$から$V^-$向き
    \end{itemize}
    に向き付けた有向グラフを$D_M$とする．
    \onslide<2->{
        また，
    \begin{itemize}
        \item $U^+ :=$ $M$に接続していない$V^+$の頂点集合
        \item $U^- :=$ $M$に接続していない$V^-$の頂点集合
    \end{itemize}
    とする．
    }

    \onslide<3->{
    \begin{block}{}
        増加道 $\longleftrightarrow$ $D_M$における$U^+$から$U^-$への有向パス
    \end{block}
    ∴ 幅優先探索により$O(m)$時間で増加道が求められる \\ \hfill {\footnotesize ($m = |E|$)}
    }
    \end{column}
    \begin{column}{.3\textwidth}
        \centering
        \tikz{
            \graph[mygraph, grow right=5em]{
            mybipartitedi;
            % matching
            {[edges={mypath}] {1,2,3,4} <-> {b,e,d,f}; }
        };
        \begin{scope}[on background layer]
        \onslide<2->{
            \node[mysubset, fit=(5) (6), "$U^+$" left ]{} ;
            \node[mysubset, fit=(a), "$U^-$" right]{} ;
            \node[mysubset, fit=(c), "$U^-$" right]{} ;
        }
        \end{scope}
        \node[above=10pt of 1]{$V^+$};
        \node[above=10pt of a]{$V^-$};
        }
        $D_M$
    \end{column}
    \end{columns}
\end{frame}

\begin{frame}
    \frametitle{逐次最短路法}

    \small
    \begin{columns}[T]
    \begin{column}{.7\textwidth}
    マッチング$M$に対し，$D_M$の枝長$\ell_M$を
    \[
        \ell_M(a) := \begin{cases}
            \phantom{-}w(e)  & (a=\text{$V^+$から$V^-$向きの$e \in E$}) \\
            -w(e) & (a = \text{$V^-$から$V^+$向きの$e \in E$})
        \end{cases}
    \]
    と定める．

    \bigskip
    \onslide<2->{
    \structure{\coloremoji{🤔}（なぜこう定めるのか？）}

    交互道 or 交互サイクル$P$に関して反転操作を行った後のマッチング$M \triangle P$の重みが
    \begin{block}{}
    \[
        w(M \triangle P) = w(M) + \ell_M(P)
    \]
    \end{block}
    と書けて便利なので．
    }
    
    \end{column}
    \begin{column}{.3\textwidth}
        \centering
        \tikz{
            \graph[mygraph, grow right=5em]{
            mybipartitedi;
            % matching
            {[edges={mypath}] {1,2,3,4} <-> {b,e,d,f}; }
        };
        \begin{scope}[on background layer]
            \node[mysubset, fit=(5) (6), "$U^+$" left ]{} ;
            \node[mysubset, fit=(a), "$U^-$" right]{} ;
            \node[mysubset, fit=(c), "$U^-$" right]{} ;
        \end{scope}
        \node[above=10pt of 1]{$V^+$};
        \node[above=10pt of a]{$V^-$};
        }
        $D_M$
    \end{column}
    \end{columns}
\end{frame}

\begin{frame}{逐次最短路法}
    \small
    \highlight[UniBlue]{$w(M \triangle P) = w(M) + \ell_M(P)$}より，$\ell_M(P)$最小の増加道$P$を選ぶのが良さそう

    \begin{block}{逐次最短路法}
        \begin{algorithmic}[1]
            \STATE $M \gets \emptyset$ \COMMENT{初期マッチング}
            \WHILE{$M$に対する増加道が存在する}
                \STATE 補助ネットワーク$(D_M, \ell_M)$上で$U^+$--$U^-$最短路$P$を求める \\ \COMMENT{Bellman--Ford法で$O(mn)$時間}
                \STATE $M \gets M \triangle P$ \COMMENT{マッチング更新}
            \ENDWHILE
        \end{algorithmic}
    \end{block}
        {\footnotesize $M \triangle P$: $M$と$P$の対称差（$P$に沿って$M$の枝を反転）}

        \pause
        \bigskip
        \begin{theorem}
            逐次最短路法で求まるマッチングを$\emptyset = M_0, M_1, \cdots,M_{\mu}$とする．このとき，各$k=0, 1, \dots, \mu$に対し，$M_k$は大きさ$k$のマッチングのうち重み最小である．また，大きさ$\mu+1$以上のマッチングは存在しない．
        \end{theorem}
\end{frame}

\begin{frame}[t]{逐次最短路法の証明}
    \small
    \begin{theorem}
        逐次最短路法で求まるマッチングを$\emptyset = M_0, M_1, \cdots,M_{\mu}$とする．このとき，各$k=0, 1, \dots, \mu$に対し，$M_k$は大きさ$k$のマッチングのうち重み最小である．
    \end{theorem}

\structure{(証明)} $k$に関する帰納法．$k=0$は自明． 

$M$が大きさ$k$の最小重みマッチングで，少なくとも1つ増加道をもつと仮定する．

\bigskip
\onslide<+->
\structure{Claim.} $(D_M, \ell_M)$は$U^+$--$U^-$最短路をもつ \\
\only<+>{{\footnotesize \structure{(∵)} 負閉路が存在しないことを示せば良い．負閉路$C$が存在したと仮定する．$C$は交互サイクルであるとしてよい．すると，$M \triangle C$は大きさ$k$のマッチングで，
\[ w(M \triangle C) = w(M) + \ell_M(C) < w(M) \]となり，矛盾．\qed}}

\bigskip
\onslide<+->
$P$を$U^+$--$U^-$最短路とし，$M^+ := M \triangle P$を更新後のマッチングとする．

\bigskip
\structure{Claim.} 任意の大きさ$k+1$のマッチング$N$に対し，$w(M^+) \leq w(N)$
\end{frame}

\begin{frame}[t]{逐次最短路法の証明}
    \small
    {\footnotesize 
    \postit{$M$: 大きさ$k$の最小重みマッチング \\ 
    $P =$ $(D_M, \ell_M)$上の$U^+$--$U^-$最短路 \\ 
    $M^+ := M \triangle P$: 更新後のマッチング}
    }
\newcommand{\M}{{\textcolor{AlertOrange}{M}}}
\newcommand{\Mp}{{\textcolor{AlertOrange}{M^+}}}
\newcommand{\N}{{\textcolor{UniGreen}{N}}}

\structure{Claim.} 任意の大きさ$k+1$のマッチング$\N$に対し，$w(\Mp) \leq w(\N)$

\bigskip
\structure{(∵)} $\M$と$\N$の対称差$\M \triangle \N$は，交互道と偶サイクルからなる．いま，$|\M| < |\N|$より，この中に$\M$増加道$Q$が存在する．

\begin{center}
    \tikz[baseline=(2.base)]{\graph[mygraph,clockwise]{subgraph I_n[n=4, radius=.7cm]; 
        {[edges={mypath }] 1 -- 2; 3 -- 4; };
        {[edges={mypath2}] 2 -- 3; 4 -- 1; };
    }}
    \quad
    \begin{tabular}{c}
    \tikz{\graph[mygraph]{
        1 --[mypath] 2 --[mypath2] 3 --[mypath] 4 --[mypath2] 5;
    }} \\
    \tikz{\graph[mygraph]{
        1 --[mypath] 2 --[mypath2] 3 --[mypath] 4;
    }} \\
    \tikz{\graph[mygraph]{
        1 --[mypath2] 2 --[mypath] 3 --[mypath2] 4 --[mypath] 5 --[mypath2] 6;
    }} $Q$
    \end{tabular}
\end{center}
\pause
$P$は最短路なので，$\ell_\M(P) \leq \ell_\M(Q)$.

\bigskip
また，$\N \triangle Q$は大きさ$k$のマッチングなので，
\[
   w(\M) \leq w(\N \triangle Q) = w(\N) - \ell_\M(Q).
\]
これらより，
\[
    w(\Mp) = w(\M) + \ell_\M(P) \leq w(\N) - \ell_\M(Q) + \ell_\M(P) \leq w(\N). \qed
\]
\end{frame}

\begin{frame}{逐次最短路法の計算量}

    \structure{計算量解析}
    \begin{itemize}
        \item $M$の更新は$\mu$回 
        {\footnotesize \hfill ($\mu =$ 最大マッチングの大きさ)}
        \item 各更新は1回の$(D_M, \ell_M)$上の最短路問題 \\
        \hfill $\longrightarrow$ Bellman--Ford法で$O(mn)$時間
    \end{itemize}
    {\footnotesize \hfill ($n = |V|, m = |E|$)}
    
    \begin{theorem}[伊理~(1960)]
        逐次最短路法は，$O(mn\mu)$時間で，$k=0,1,\dots,\mu$に対する大きさ$k$の最小重みマッチング$M_k$を求める．
    \end{theorem}

    \bigskip
    Cf. ハンガリー法 $O(mn^3)$時間
\end{frame}

\begin{frame}{主双対法}
    \small
    マッチング$M$とポテンシャル$p$の両方を保持することで，逐次最短路法の計算量を改善できる．

    \begin{definition}[簡約枝長]
        枝長$\ell: A \to \R$のポテンシャル$p: V \to \R$に関する\alert{簡約枝長}$\ell_p: A \to \R$を次のように定める:
        \begin{align*}
       \ell_p(a) := \ell(a) + p(i) - p(j) \quad (a = ij \in A) 
        \end{align*}
    \end{definition}

    \bigskip
    \structure{観察}
    \begin{itemize}
        \item 
    任意の$s$--$t$路$P$に対し，
    \[
        \ell_p(P) = \ell(P) + p(s) - p(t).
    \]
    特に，$\ell_p$に関する$s$--$t$最短路 $\iff$ $\ell$に関する$s$--$t$最短路．
        \item $\ell_p \geq 0$ \dotfill\alert{Bellman-Ford法ではなくDijkstra法が使える！}
    \end{itemize}
\end{frame}

\begin{frame}{主双対法}
    
    \small
    \begin{block}{主双対法}
        \begin{algorithmic}[1]
            \STATE $M := \emptyset$ \COMMENT{初期マッチング}
            \STATE $p(i) := 0$ ($i \in V^+$), $p(j) := \min_{i \in \delta(i)} w(ij)$ ($j \in V^-$) \COMMENT{初期ポテンシャル}
            \WHILE{$M$に対する増加道が存在する}
                \STATE 簡約枝長$\ell_{M,p}$に関する$D_M$上の$U^+$--$U^-$最短路$P$と最短路長関数$d$を求める． 
                \COMMENT{Dijkstra法で$O(m + n\log n)$時間}
                \STATE $M \gets M \triangle P$, $p \gets p + d$ \COMMENT{マッチングとポテンシャル更新}
            \ENDWHILE
        \end{algorithmic} 
    \end{block}

    \vfill
    \begin{theorem}[Edmonds--Karp~(1970), 冨澤~(1971)]
        主双対法で求まるマッチングを$\emptyset = M_0, M_1, \cdots,M_{\mu}$とする（$\mu =$最大マッチングの大きさ）．このとき，各$k=0, 1, \dots, \mu$に対し，$M_k$は大きさ$k$のマッチングのうち重み最小である．また，全体の計算量は$O((m + n \log n)\mu)$時間．
    \end{theorem}
\end{frame}

\begin{frame}{主双対法の証明}
    \small
    次の補題を示せば，あとは逐次最短路法の定理から従う．
    
    \begin{lemma}
    主双対法における$p$は，常に$(D_M, \ell_M)$上のポテンシャル．
    \end{lemma}
    \pause
    \structure{(証明)} 
    $p^+ := p + d$とする．{\footnotesize ($d: V \to \R$ ... $(D_M, \ell_{M,p})$上の（$U^+$からの）最短路長関数)}

    \bigskip
    $d$は$(D_M, \ell_{M,p})$のポテンシャルなので，任意の枝$a = ij$に対し，
    \begin{align*}
        d(j) - d(i) \leq \ell_{M,p}(a) = \ell_M(a) + p(i) - p(j)
    \end{align*}
    が成り立つ．移項すると $p^+(j) - p^+(i) \leq \ell_M(a)$ となり，$p^+$は$(D_M, \ell_M)$のポテンシャルである．

    \pause
    \bigskip
    次に，$M \gets M \triangle P$と更新した後も$p^+$がポテンシャルであることを示す．
    最短路$P$上の枝$a = ij$では，上の不等式で等号が成立:
    \begin{align*}
        d(j) - d(i) = \ell_{M,p}(a) \quad \iff \quad p^+(j) - p^+(i) = \ell_M(a)
    \end{align*}
    よって，$P$の枝を反転しても，$p^+$がポテンシャルであることは保たれる． \qed
\end{frame}

%%%%%%%%%%%%%%%%%%%%%%%%%%%%%%%%%%%%%%%%%%%%%
\section{最適性条件からの見方}
\againframe<4|handout:0>{outline}
%%%%%%%%%%%%%%%%%%%%%%%%%%%%%%%%%%%%%%%%%%%%%

\begin{frame}{最適性条件からの見方}
    ハンガリー法，逐次最短路法，主双対法は，どれも\textbf{“同じアルゴリズム”の異なる姿}と思える．

    \bigskip\pause
    以下，最小重み\textbf{完全}マッチングを考えよう．

    \vfill
    \small
    マッチング$M$が\alert{完全} $\defiff$ 全ての頂点が$M$に接続 $\iff$ $2|M| = |V|$

    \begin{block}{重み付き二部完全マッチング問題}
        \structure{入力} $G=(V; E)$: 二部グラフ($|V^+|=|V^-|=n/2$)，枝重み$w_e$ ($e \in E$) \\
        \structure{出力} $w(M) := \sum_{e\in M}w_e$が最小の完全マッチング$M$
    \end{block}
    ※ 以下，$G$は少なくとも1つの完全マッチングを持つとする．
\end{frame}

\begin{frame}{双対変数とポテンシャル}
    \small
    \begin{columns}[T]
    \begin{column}{.45\textwidth}
    \begin{primal}
        \setlength{\abovedisplayskip}{0pt}
        \begin{alignat*}{2}
            \text{min} & \quad  \sum_{e \in E} w_e x_e \\
            \text{s.t.} & \quad \sum_{e \in \delta(i)} x_e = 1  \quad (i \in V) \\
            & \quad x \geq 0 
        \end{alignat*}
    \end{primal}
    \end{column}
    \begin{column}{.45\textwidth}
    \begin{dual}
        \setlength{\abovedisplayskip}{0pt}
        \begin{alignat*}{2}
            \text{max} & \quad  \sum_{i \in V} y_i =: y(V) \\
            \text{s.t.} & \quad y_i + y_j \leq w_e \quad (e=ij \in E) 
        \end{alignat*}
    \end{dual}
    \centering
    \begin{block}{}
        \centering
    双対変数$y \in \R^V$ s.t. \\ $y_i + y_j \leq w_e$ ($e = ij \in E$) 
    \end{block}
    $\updownarrow$ \\
    \begin{alertblock}{}
        \centering
    ポテンシャル$p: V \to \R$ s.t. \\ $p(j) - p(i) \leq w_e$ ($e = ij \in E$) 
    \end{alertblock}
    \end{column}
    \end{columns}
\end{frame}

\begin{frame}{最適性条件のポテンシャルを用いた書き換え}
    
    \begin{columns}[T]
    \begin{column}{.6\textwidth}
        \begin{block}{最適性条件 (復習)}
        $M \subseteq E$, $y \in \R^V$が最適解 
        $\iff$
        \begin{enumerate}
            \item $M$は完全マッチング
            \item $y$は(D)の実行可能解
            \item $y_i + y_j = w_e$ \quad ($e = ij \in M$)
        \end{enumerate}
        \end{block}

        \begin{alertblock}{補題}
            $M\subseteq E$, $p: V \to \R$が最適解
            $\iff$
            \begin{enumerate}
                \item $M$は完全マッチング
                \item $p$は$(D_M, \ell_M)$のポテンシャル
            \end{enumerate}
        \end{alertblock}
    \end{column}
    \begin{column}{.35\textwidth}
        \footnotesize
        \begin{primal}
        \begin{alignat*}{2}
            \text{min} & \quad  \sum_{e \in E} w_e x_e \\
            \text{s.t.} & \quad \sum_{e \in \delta(i)} x_e = 1  \quad (i \in V) \\
            & \quad x \geq 0 
        \end{alignat*}
    \end{primal}
    \begin{dual}
        \setlength{\abovedisplayskip}{0pt}
        \begin{alignat*}{2}
            \text{max} & \quad  \sum_{i \in V} y_i =: y(V) \\
            \text{s.t.} & \quad y_i + y_j \leq w_e \quad (e \in E) 
        \end{alignat*}
    \end{dual}
    \end{column}
    \end{columns}
\end{frame}

\begin{frame}{最適性条件のポテンシャルを用いた書き換え}
    \small
    \begin{columns}[T]
    \begin{column}{.48\textwidth}
        \begin{block}{最適性条件 (復習)}
        $M \subseteq E$, $y \in \R^V$が最適解 
        $\iff$
        \begin{enumerate}
            \setlength{\itemsep}{-0.3em}
            \item $M$は完全マッチング
            \item $y$は(D)の実行可能解
            \item $y_i + y_j = w_e$ \quad ($e = ij \in M$)
        \end{enumerate}
        \end{block}
    \end{column}
    \begin{column}{.48\textwidth}
        \begin{alertblock}{補題}
            $M\subseteq E$, $p: V \to \R$が最適解
            $\iff$
            \begin{enumerate}
                \setlength{\itemsep}{-0.3em}
                \item $M$は完全マッチング
                \item $p$は$(D_M, \ell_M)$のポテンシャル
            \end{enumerate}
        \end{alertblock}
    \end{column}
    \end{columns}

    \bigskip
    \structure{(証明)}
    $D_M$においては，$M$の枝は両向きであることに注意する．
    $p$がポテンシャルならば，$e = ij \in M$に対し，
    \begin{align*}
        p(j) - p(i) & \leq \ell_M(\overrightarrow{e}) = w_e \\
        p(i) - p(j) & \leq \ell_M(\overleftarrow{e}) = -w_e 
    \end{align*}
    が成り立つので，対応する$y$は\structure{③}を満たす．逆も同様．\qed
\end{frame}

\begin{frame}{ポテンシャル最適性条件・閉路最適性条件}
    \begin{lemma}
        完全マッチング$M$に対し，以下は同値:
        \begin{enumerate}
            \item $M$は最小重み完全マッチング
            \item $(D_M, \ell_M)$のポテンシャル$p$が存在
            \item $(D_M, \ell_M)$に負閉路は存在しない
        \end{enumerate}
    \end{lemma}

    \small
    \bigskip
    \structure{(証明)} \\
    \textbf{①$\iff$②:} $M$が最適 $\iff$ $\exists p$ s.t. $(M, p)$は相補性条件を満たす \\
    \textbf{②$\iff$③:} 最短路問題のところでやった． \qed
\end{frame}


\begin{frame}{ハンガリー法（ポテンシャル版）}
    \small
    \begin{block}{ハンガリー法}
        \begin{algorithmic}[1]
            \STATE $p(i) := 0$ ($i \in V^+$), $p(j) := \min_{e \in \delta(j)} w_e$ ($j \in V^-$)とする．
            \COMMENT{$p$は初期ポテンシャル}
            \WHILE{\textbf{True}}
                \onslide<1>{\STATE $G_p$の（重みなし）最大マッチング$M$を求める． \COMMENT{$O(mn)$時間}} %
                \IF{\alt<1>{$M$が完全マッチング}{$G_p$が完全マッチング(say $M$)をもつ}}
                    \RETURN $M$ 
                    \COMMENT{相補性条件より，$M$は最小重みマッチング}
                \ELSE[{ポテンシャル$p$を更新}]
                    \STATE \alt<1>{$X$を$D_M(p)$における$U^+$から到達可能な頂点全体とする．}{$G_p$の最小頂点被覆$(S,T)$を求める \COMMENT{$|S|+|T| < n$}}
                    \STATE \alt<1>{$\eps := \min\{ w_{e} + p(i) - p(j) : e = ij, i \in V^+ \cap X, j \in V^- \setminus X \}$}{$\eps := \min\{ w_{e} + p(i) - p(j) : e = ij \in E \text{ s.t. $(S, T)$に被覆されていない} \}$}
                    \STATE $p(i) \gets \begin{cases}
                        p(i) - \eps & (i \in \alt<1>{X}{(V^+ \setminus S) \cup T}) \\
                        p(i)  & (\text{otherwise})
                    \end{cases}$
                \ENDIF
            \ENDWHILE
        \end{algorithmic} 
    \end{block}

    \bigskip
    \onslide<2->
    最後の反復以外，$M$を使っていない！実質的に保持しているのはポテンシャル$p$のみ．
\end{frame}

\begin{frame}{逐次最短路法，主双対法，ハンガリー法}

    {\small
    \begin{center}
        \renewcommand{\arraystretch}{1.7}
        \begin{NiceTabular}{c|ccc}
            \CodeBefore
                \rowcolors{1}{UniBlue!10}{}
            \Body
            & \bf 逐次最短路法 & \bf 主双対法 & \bf ハンガリー法 \\\hline
            変数 & \Block[]{}{マッチング$M$ \\ \footnotesize (ポテンシャルは陰に保持)} & \Block[]{}{マッチング$M$ \\ ポテンシャル$p$} & \Block[]{}{\footnotesize (マッチングは保持しない) \\ ポテンシャル$p$} \\
            更新方法 & 最短路(Bellman--Ford) & 最短路(Dijkstra) & 幅優先探索 \\
            計算量 & $O(mn^2)$ & $O((m + n \log n)n)$ & $O(mn^3)$ \\
        \end{NiceTabular}
    \end{center}
    }

    \bigskip
    どれも
    \begin{block}{}
        \centering
    ポテンシャル（双対変数）$p:V \to \R$を保持し，\\ マッチング$M$が完全になるまで更新する
    \end{block}
    アルゴリズムとみなせる．
\end{frame}


%%%%%%%%%%%%%%%%%%%%%%%%%%%%%%%%%%%%%%%%%%%%%%
\section{まとめ}
\begin{frame}{まとめ: 二部マッチング}

    \structure{まとめ}
    \begin{itemize}
        \item 重みなし: 増加道アルゴリズム
        \item 重みあり: 逐次最短路法，主双対法，ハンガリー法
        \item 双対変数と最短路問題のポテンシャルの関係
    \end{itemize}
    \dotfill ネットワークフローやマトロイド交差へ

    \bigskip
    \structure{その他のアルゴリズム}
    \begin{itemize}
        \item Hopcroft--Karp--Karzanov法 (1973): 重みなし$O(m\sqrt{n})$時間 \\
        \item 重みあり: $m^{1 + o(1)}$時間（\emph{JACM} 2025\footnote{L. Chen, R. Kyng, Y. P. Liu, R. Peng, M. P. Gutenberg, and S. Sachdeva, \textit{``Maximum Flow and Minimum-Cost Flow in Almost-Linear Time''}, \textit{JACM}, 2025.}, 最小費用流）
    \end{itemize}
    
\end{frame}
\end{document}
